% --- Archon physics formalization source begin ---
% archon:physics
% archon:covers PhyXMiniProblems/problem_phyx_mini_0049.lean
% archon:source-report reports/phyx_mini/problem_phyx_mini_0049.source.json

\chapter{Physics problem phyx\_mini\_0049}
\label{ch:PhyXMiniProblems_problem_phyx_mini_0049}

\paragraph{Problem source.}
\#\# Physical scenario

A laser beam is aimed at a 1.0\textbackslash{}text\{-\}\textbackslash{}mathrm\{cm\}\textbackslash{}text\{-thick\} sheet of glass at an angle 30\textasciicircum{}\textbackslash{}circ above the glass.

\#\# Question

What is its direction in the air on the other side?

\#\# Answer choices

- A: A: 64.0\textasciicircum{}\textbackslash{}circ

- B: B: 60.0\textasciicircum{}\textbackslash{}circ

- C: C: 50.0\textasciicircum{}\textbackslash{}circ

- D: D: 56.0\textasciicircum{}\textbackslash{}circ

\#\# Figure caption (auxiliary; use the image as primary evidence)

The image is a geometric diagram showing a room with dimensions and a light bulb. Here's a breakdown of the content:

- **Room Dimensions**: The room has a height of 2.50 meters and a width of 1.00 meter along the ceiling, with the distance from the bottom-left corner to the right side being represented as \textbackslash{}( l\_1 + l\_2 \textbackslash{}).

- **Light Bulb**: Positioned at a corner near the top right of the room, labeled as "Bulb."

- **Mirror**: There is a vertical mirror located on the right side of the diagram with a height of 1.50 meters and placed 0.50 meters from the top.

- **Angles and Lines**:
  - There are lines representing the path of light emanating from the bulb, reflecting off the mirror, and reaching the bottom left corner of the room.
  - Two angles at the reflection and incidence points are marked as \textbackslash{}( \textbackslash{}theta\_1 \textbackslash{}) and \textbackslash{}( \textbackslash{}theta\_2 \textbackslash{}).

- **Distances**:
  - The ceiling distance from the wall to the bulb is labeled as 1.00 meter.
  - Various lines and angles within the diagram illustrate the interaction of light from the bulb with the mirror.

- **Text Labels**:
  - The diagram includes measurements for height and distance.
  - \textbackslash{}( \textbackslash{}theta\_1 \textbackslash{}) and \textbackslash{}( \textbackslash{}theta\_2 \textbackslash{}) denote angles involving the light’s path.

This setup likely illustrates the geometric principles of light reflection and the law of reflection.

\#\# Dataset metadata

Geometrical Optics; Spatial Relation Reasoning, Physical Model Grounding Reasoning

\paragraph{Recorded answer/context.}
B: B: 60.0\textasciicircum{}\textbackslash{}circ

\paragraph{Figure/image path.}
/root/proposal\_for\_physic/science-mango/phyx\_mini\_run/phyx\_data/test\_image/49.png

\paragraph{Formalization target.}
create a compiling Lean file with sorry bodies at `PhyXMiniProblems/problem\_phyx\_mini\_0049.lean`. The Lean declarations must preserve the physical quantities, dimensions or dimensional roles, figure labels, governing-law hypotheses, and final relation expressed by this problem.
Use Mathlib/Physlib names found through LeanExplore where available. If a physics API is missing, introduce faithful local abstractions rather than scalar placeholder aliases.

\begin{theorem}[Physics formalization target]
\label{thm:physics:phyx_mini_0049:target}
\lean{PhyXMiniProblems.ProblemPhyXMini0049.outgoingDirectionIsAnswerB}
\uses{def:physics:phyx-mini-0049:phyxminiproblems-problemphyxmini0049-degreestoradians, def:physics:phyx-mini-0049:phyxminiproblems-problemphyxmini0049-parallelglasssheetdiagram, def:physics:phyx-mini-0049:phyxminiproblems-problemphyxmini0049-matchesglasssheetproblemdata, def:physics:phyx-mini-0049:phyxminiproblems-problemphyxmini0049-hasphysicalopticalparameters, def:physics:phyx-mini-0049:phyxminiproblems-problemphyxmini0049-satisfiesparallelsheetgeometry, def:physics:phyx-mini-0049:phyxminiproblems-problemphyxmini0049-obeyssnellslaw, def:physics:phyx-mini-0049:phyxminiproblems-problemphyxmini0049-answerchoice}
The assigned autoformalize agent should translate this physics problem into Lean declarations in the covered file, with theorem and lemma proof bodies written as `by sorry`.
\end{theorem}
\begin{proof}
This is an autoformalization task, not a proof task. Produce statements that can later be proved by the physics prover without weakening the physical model.
\end{proof}

% --- Archon declaration topology begin ---
\section{Declaration topology}
\begin{definition}[Dim Length]
  \label{def:physics:phyx-mini-0049:phyxminiproblems-problemphyxmini0049-dimlength}
  \lean{PhyXMiniProblems.ProblemPhyXMini0049.DimLength}
  A physical length whose scalar carrier is real-valued
\end{definition}
\begin{proof}
  This is the declared model definition of Dim Length.
\end{proof}
\begin{definition}[centimeter Unit Choices]
  \label{def:physics:phyx-mini-0049:phyxminiproblems-problemphyxmini0049-centimeterunitchoices}
  \lean{PhyXMiniProblems.ProblemPhyXMini0049.centimeterUnitChoices}
  Unit choices in which length readouts are expressed in centimeters
\end{definition}
\begin{proof}
  This is the declared model definition of centimeter Unit Choices.
\end{proof}
\begin{definition}[one Centimeter]
  \label{def:physics:phyx-mini-0049:phyxminiproblems-problemphyxmini0049-onecentimeter}
  \lean{PhyXMiniProblems.ProblemPhyXMini0049.oneCentimeter}
  \uses{def:physics:phyx-mini-0049:phyxminiproblems-problemphyxmini0049-dimlength, def:physics:phyx-mini-0049:phyxminiproblems-problemphyxmini0049-centimeterunitchoices}
  The dimensionful sheet thickness specified by the problem, 1.0 cm
\end{definition}
\begin{proof}
  This is the declared model definition of one Centimeter.
\end{proof}
\begin{definition}[degrees To Radians]
  \label{def:physics:phyx-mini-0049:phyxminiproblems-problemphyxmini0049-degreestoradians}
  \lean{PhyXMiniProblems.ProblemPhyXMini0049.degreesToRadians}
  Convert a degree readout to the radian scalar used by Real.sin
\end{definition}
\begin{proof}
  This is the declared model definition of degrees To Radians.
\end{proof}
\begin{definition}[Diagram Plane]
  \label{def:physics:phyx-mini-0049:phyxminiproblems-problemphyxmini0049-diagramplane}
  \lean{PhyXMiniProblems.ProblemPhyXMini0049.DiagramPlane}
  The two-dimensional plane of the ray diagram
\end{definition}
\begin{proof}
  This is the declared model definition of Diagram Plane.
\end{proof}
\begin{definition}[Ray Direction]
  \label{def:physics:phyx-mini-0049:phyxminiproblems-problemphyxmini0049-raydirection}
  \lean{PhyXMiniProblems.ProblemPhyXMini0049.RayDirection}
  \uses{def:physics:phyx-mini-0049:phyxminiproblems-problemphyxmini0049-diagramplane}
  A geometric propagation direction is a nonzero vector in the diagram plane
\end{definition}
\begin{proof}
  This is the declared model definition of Ray Direction.
\end{proof}
\begin{definition}[Optical Region]
  \label{def:physics:phyx-mini-0049:phyxminiproblems-problemphyxmini0049-opticalregion}
  \lean{PhyXMiniProblems.ProblemPhyXMini0049.OpticalRegion}
  The three homogeneous optical regions traversed in propagation order
\end{definition}
\begin{proof}
  This is the declared model definition of Optical Region.
\end{proof}
\begin{definition}[Sheet Interface]
  \label{def:physics:phyx-mini-0049:phyxminiproblems-problemphyxmini0049-sheetinterface}
  \lean{PhyXMiniProblems.ProblemPhyXMini0049.SheetInterface}
  The parallel entry and exit faces of the sheet
\end{definition}
\begin{proof}
  This is the declared model definition of Sheet Interface.
\end{proof}
\begin{definition}[Ray Segment]
  \label{def:physics:phyx-mini-0049:phyxminiproblems-problemphyxmini0049-raysegment}
  \lean{PhyXMiniProblems.ProblemPhyXMini0049.RaySegment}
  The three directed portions of the laser path
\end{definition}
\begin{proof}
  This is the declared model definition of Ray Segment.
\end{proof}
\begin{definition}[incident Region]
  \label{def:physics:phyx-mini-0049:phyxminiproblems-problemphyxmini0049-incidentregion}
  \lean{PhyXMiniProblems.ProblemPhyXMini0049.incidentRegion}
  \uses{def:physics:phyx-mini-0049:phyxminiproblems-problemphyxmini0049-opticalregion, def:physics:phyx-mini-0049:phyxminiproblems-problemphyxmini0049-sheetinterface}
  Incident optical region at each interface, in propagation order
\end{definition}
\begin{proof}
  This is the declared model definition of incident Region.
\end{proof}
\begin{definition}[transmitted Region]
  \label{def:physics:phyx-mini-0049:phyxminiproblems-problemphyxmini0049-transmittedregion}
  \lean{PhyXMiniProblems.ProblemPhyXMini0049.transmittedRegion}
  \uses{def:physics:phyx-mini-0049:phyxminiproblems-problemphyxmini0049-opticalregion, def:physics:phyx-mini-0049:phyxminiproblems-problemphyxmini0049-sheetinterface}
  Transmitted optical region at each interface, in propagation order
\end{definition}
\begin{proof}
  This is the declared model definition of transmitted Region.
\end{proof}
\begin{definition}[incident Segment]
  \label{def:physics:phyx-mini-0049:phyxminiproblems-problemphyxmini0049-incidentsegment}
  \lean{PhyXMiniProblems.ProblemPhyXMini0049.incidentSegment}
  \uses{def:physics:phyx-mini-0049:phyxminiproblems-problemphyxmini0049-sheetinterface, def:physics:phyx-mini-0049:phyxminiproblems-problemphyxmini0049-raysegment}
  Ray segment incident on each interface
\end{definition}
\begin{proof}
  This is the declared model definition of incident Segment.
\end{proof}
\begin{definition}[transmitted Segment]
  \label{def:physics:phyx-mini-0049:phyxminiproblems-problemphyxmini0049-transmittedsegment}
  \lean{PhyXMiniProblems.ProblemPhyXMini0049.transmittedSegment}
  \uses{def:physics:phyx-mini-0049:phyxminiproblems-problemphyxmini0049-sheetinterface, def:physics:phyx-mini-0049:phyxminiproblems-problemphyxmini0049-raysegment}
  Ray segment transmitted through each interface
\end{definition}
\begin{proof}
  This is the declared model definition of transmitted Segment.
\end{proof}
\begin{definition}[Parallel Glass Sheet Diagram]
  \label{def:physics:phyx-mini-0049:phyxminiproblems-problemphyxmini0049-parallelglasssheetdiagram}
  \lean{PhyXMiniProblems.ProblemPhyXMini0049.ParallelGlassSheetDiagram}
  \uses{def:physics:phyx-mini-0049:phyxminiproblems-problemphyxmini0049-dimlength, def:physics:phyx-mini-0049:phyxminiproblems-problemphyxmini0049-raydirection, def:physics:phyx-mini-0049:phyxminiproblems-problemphyxmini0049-opticalregion, def:physics:phyx-mini-0049:phyxminiproblems-problemphyxmini0049-sheetinterface, def:physics:phyx-mini-0049:phyxminiproblems-problemphyxmini0049-raysegment}
  Physical quantities and labelled geometry for the parallel glass sheet. Each face normal points in the direction of propagation through that face. The incidence and refraction angles are scalar radian readouts measured from those normals. incomingElevationRadians is instead measured above the sheet surface, matching the wording of the problem
\end{definition}
\begin{proof}
  This is the declared model definition of Parallel Glass Sheet Diagram.
\end{proof}
\begin{definition}[angle Between Directions]
  \label{def:physics:phyx-mini-0049:phyxminiproblems-problemphyxmini0049-anglebetweendirections}
  \lean{PhyXMiniProblems.ProblemPhyXMini0049.angleBetweenDirections}
  \uses{def:physics:phyx-mini-0049:phyxminiproblems-problemphyxmini0049-raydirection}
  The undirected radian angle between two nonzero diagram directions
\end{definition}
\begin{proof}
  This is the declared model definition of angle Between Directions.
\end{proof}
\begin{definition}[Is Physical Acute Angle]
  \label{def:physics:phyx-mini-0049:phyxminiproblems-problemphyxmini0049-isphysicalacuteangle}
  \lean{PhyXMiniProblems.ProblemPhyXMini0049.IsPhysicalAcuteAngle}
  An angle readout lies on the acute geometrical-optics branch
\end{definition}
\begin{proof}
  This is the declared model definition of Is Physical Acute Angle.
\end{proof}
\begin{definition}[Matches Glass Sheet Problem Data]
  \label{def:physics:phyx-mini-0049:phyxminiproblems-problemphyxmini0049-matchesglasssheetproblemdata}
  \lean{PhyXMiniProblems.ProblemPhyXMini0049.MatchesGlassSheetProblemData}
  \uses{def:physics:phyx-mini-0049:phyxminiproblems-problemphyxmini0049-onecentimeter, def:physics:phyx-mini-0049:phyxminiproblems-problemphyxmini0049-degreestoradians, def:physics:phyx-mini-0049:phyxminiproblems-problemphyxmini0049-parallelglasssheetdiagram}
  Textual problem data: a 1.0 cm glass sheet, a 30° incoming elevation, and the same ambient air on both sides. No outgoing-angle value is recorded here
\end{definition}
\begin{proof}
  This is the declared model definition of Matches Glass Sheet Problem Data.
\end{proof}
\begin{definition}[Has Physical Optical Parameters]
  \label{def:physics:phyx-mini-0049:phyxminiproblems-problemphyxmini0049-hasphysicalopticalparameters}
  \lean{PhyXMiniProblems.ProblemPhyXMini0049.HasPhysicalOpticalParameters}
  \uses{def:physics:phyx-mini-0049:phyxminiproblems-problemphyxmini0049-parallelglasssheetdiagram, def:physics:phyx-mini-0049:phyxminiproblems-problemphyxmini0049-isphysicalacuteangle}
  Positivity and principal-branch conditions for the optical model
\end{definition}
\begin{proof}
  This is the declared model definition of Has Physical Optical Parameters.
\end{proof}
\begin{definition}[Satisfies Parallel Sheet Geometry]
  \label{def:physics:phyx-mini-0049:phyxminiproblems-problemphyxmini0049-satisfiesparallelsheetgeometry}
  \lean{PhyXMiniProblems.ProblemPhyXMini0049.SatisfiesParallelSheetGeometry}
  \uses{def:physics:phyx-mini-0049:phyxminiproblems-problemphyxmini0049-incidentsegment, def:physics:phyx-mini-0049:phyxminiproblems-problemphyxmini0049-transmittedsegment, def:physics:phyx-mini-0049:phyxminiproblems-problemphyxmini0049-parallelglasssheetdiagram, def:physics:phyx-mini-0049:phyxminiproblems-problemphyxmini0049-anglebetweendirections}
  Geometric relations supplied by the ray diagram and the parallel planar faces. In particular, the internal ray meets both equal-oriented normals at the same angle. The incoming elevation and entry incidence are complementary
\end{definition}
\begin{proof}
  This is the declared model definition of Satisfies Parallel Sheet Geometry.
\end{proof}
\begin{definition}[Satisfies Snells Law At]
  \label{def:physics:phyx-mini-0049:phyxminiproblems-problemphyxmini0049-satisfiessnellslawat}
  \lean{PhyXMiniProblems.ProblemPhyXMini0049.SatisfiesSnellsLawAt}
  \uses{def:physics:phyx-mini-0049:phyxminiproblems-problemphyxmini0049-sheetinterface, def:physics:phyx-mini-0049:phyxminiproblems-problemphyxmini0049-incidentregion, def:physics:phyx-mini-0049:phyxminiproblems-problemphyxmini0049-transmittedregion, def:physics:phyx-mini-0049:phyxminiproblems-problemphyxmini0049-parallelglasssheetdiagram}
  Snell's law n₁ sin θ₁ = n₂ sin θ₂ at one sheet interface
\end{definition}
\begin{proof}
  This is the declared model definition of Satisfies Snells Law At.
\end{proof}
\begin{definition}[Obeys Snells Law]
  \label{def:physics:phyx-mini-0049:phyxminiproblems-problemphyxmini0049-obeyssnellslaw}
  \lean{PhyXMiniProblems.ProblemPhyXMini0049.ObeysSnellsLaw}
  \uses{def:physics:phyx-mini-0049:phyxminiproblems-problemphyxmini0049-parallelglasssheetdiagram, def:physics:phyx-mini-0049:phyxminiproblems-problemphyxmini0049-satisfiessnellslawat}
  The laser obeys Snell's law at both faces of the glass sheet
\end{definition}
\begin{proof}
  This is the declared model definition of Obeys Snells Law.
\end{proof}
\begin{definition}[Auxiliary Mirror Room Figure]
  \label{def:physics:phyx-mini-0049:phyxminiproblems-problemphyxmini0049-auxiliarymirrorroomfigure}
  \lean{PhyXMiniProblems.ProblemPhyXMini0049.AuxiliaryMirrorRoomFigure}
  Labels visible in the supplied auxiliary image, which depicts a bulb and a vertical mirror rather than the glass-sheet scenario. Distances are scalar meter readouts; thetaOneRadians and thetaTwoRadians are the displayed angle labels. The unknown horizontal labels l₁ and l₂ remain parameters
\end{definition}
\begin{proof}
  This is the declared model definition of Auxiliary Mirror Room Figure.
\end{proof}
\begin{definition}[Matches Auxiliary Mirror Room Image]
  \label{def:physics:phyx-mini-0049:phyxminiproblems-problemphyxmini0049-matchesauxiliarymirrorroomimage}
  \lean{PhyXMiniProblems.ProblemPhyXMini0049.MatchesAuxiliaryMirrorRoomImage}
  \uses{def:physics:phyx-mini-0049:phyxminiproblems-problemphyxmini0049-auxiliarymirrorroomfigure}
  Numerical length readouts visible in the mismatched auxiliary image
\end{definition}
\begin{proof}
  This is the declared model definition of Matches Auxiliary Mirror Room Image.
\end{proof}
\begin{definition}[Answer Choice]
  \label{def:physics:phyx-mini-0049:phyxminiproblems-problemphyxmini0049-answerchoice}
  \lean{PhyXMiniProblems.ProblemPhyXMini0049.AnswerChoice}
  The four displayed answer labels
\end{definition}
\begin{proof}
  This is the declared model definition of Answer Choice.
\end{proof}
\begin{definition}[direction Degrees]
  \label{def:physics:phyx-mini-0049:phyxminiproblems-problemphyxmini0049-answerchoice-directiondegrees}
  \lean{PhyXMiniProblems.ProblemPhyXMini0049.AnswerChoice.directionDegrees}
  \uses{def:physics:phyx-mini-0049:phyxminiproblems-problemphyxmini0049-answerchoice}
  Direction-angle readout in degrees printed beside each answer choice
\end{definition}
\begin{proof}
  This is the declared model definition of direction Degrees.
\end{proof}
\begin{definition}[recorded Answer Choice]
  \label{def:physics:phyx-mini-0049:phyxminiproblems-problemphyxmini0049-recordedanswerchoice}
  \lean{PhyXMiniProblems.ProblemPhyXMini0049.recordedAnswerChoice}
  \uses{def:physics:phyx-mini-0049:phyxminiproblems-problemphyxmini0049-answerchoice}
  Dataset metadata: the recorded answer label, not a theorem premise
\end{definition}
\begin{proof}
  This is the declared model definition of recorded Answer Choice.
\end{proof}
% --- Archon declaration topology end ---
% --- Archon physics formalization source end ---
